\documentclass[10pt,twocolumn,a4paper]{article}

\usepackage{xetexko}
\usepackage{geometry}
\usepackage{titlesec}
\usepackage{booktabs}
\usepackage{array}
\usepackage{tabularx}
\usepackage{setspace}
\usepackage{indentfirst}
\usepackage{hyperref}
\usepackage{caption}
\usepackage{float}

\geometry{
  a4paper,
  top=22mm, bottom=22mm,
  left=18mm, right=18mm,
  columnsep=6mm
}

\setmainfont{Apple SD Gothic Neo}[
  BoldFont = Apple SD Gothic Neo Bold,
  ItalicFont = Apple SD Gothic Neo Regular
]
\setsansfont{Apple SD Gothic Neo}
\setmonofont{Courier New}
\setmainhangulfont{Apple SD Gothic Neo}[
  BoldFont = Apple SD Gothic Neo Bold
]

% KCC 양식: 아라비아 숫자, 좌측 정렬, 굵게
\titleformat{\section}[block]
  {\normalsize\bfseries}
  {\thesection.\,}{0.3em}{}
\titlespacing{\section}{0pt}{8pt}{4pt}

\titleformat{\subsection}[block]
  {\normalsize\bfseries}
  {\thesubsection\,}{0.3em}{}
\titlespacing{\subsection}{0pt}{6pt}{3pt}

\setlength{\parindent}{1em}
\setlength{\parskip}{2pt}

\renewcommand{\footnoterule}{\vspace{-3pt}\noindent\rule{0.4\columnwidth}{0.4pt}\vspace{3pt}}

% 페이지 하단 가운데 페이지 번호만 표시 (헤더 없음)
\pagestyle{plain}

\begin{document}

\twocolumn[{%
\begin{center}
  {\Large\bfseries RAG와 페르소나 LLM 기반}\\[4pt]
  {\Large\bfseries 제주 설화·민담 여행 코스 추천 서비스 개발$^{*}$}\\[6pt]
  {\normalsize --- 최종 보고서 ---}\\[8pt]

  {\normalsize 조익준$^{\circ}$}\\[2pt]
  {\normalsize 경희대학교 컴퓨터공학과}\\[2pt]
  {\small\itshape \underline{keepdev0919@gmail.com}}\\[10pt]

  {\normalsize\bfseries RAG and Persona-Based LLM Service for}\\[2pt]
  {\normalsize\bfseries Jeju Folklore Travel Course Recommendation}\\[4pt]
  {\normalsize Ikjun Cho$^{\circ}$}\\[2pt]
  {\normalsize School of Computer Science and Engineering,}\\[2pt]
  {\normalsize Kyung Hee University}\\[12pt]
\end{center}
\noindent\rule{\textwidth}{0.4pt}
\vspace{4pt}
}]

% 첫 페이지 하단 각주 - 연구 지원 표시
\thispagestyle{plain}
\renewcommand{\thefootnote}{*}
\footnotetext{본 연구는 경희대학교 컴퓨터공학과 졸업프로젝트로 수행되었으며,
본 결과물은 2026 관광데이터 활용 공모전(예비 합격) 출품 예정 작품이다.}
\renewcommand{\thefootnote}{\arabic{footnote}}
\setcounter{footnote}{0}

% ── 요약 ─────────────────────────────────────────────────
\begin{center}
{\normalsize\bfseries 요~~~약}
\end{center}

{\small\noindent
본 연구는 제주도 방문 관광객을 대상으로, 실제 관광객이 검증한 여행 코스에
LLM이 설화·민담 콘텐츠를 매핑하여 \emph{계획·탐험·기록}이 하나의 여정으로
연결되는 iOS 모바일 서비스를 구현한다. LangGraph 인터페이스를 모방한 2단계
LLM 함수 파이프라인이 지역·스타일·기간 입력을 받아 코스 후보 풀에서 무작위 3건을
추천하고, 5종 페르소나 기반 AI 설화 동행자가 현장에서 RAG 기반 설화 해설을
SSE(Server-Sent Events) 스트리밍과 페르소나 전용 TTS 음성으로 동시 제공한다.
여행 후 GPT-4o가 방문 이력·대화 맥락·커뮤니티 리뷰를 종합하여 1인칭 일지와
한국 민화 스타일 이미지를 단일 호출의 구조화 출력으로 동시 생성한다.
데이터 파이프라인은 제주도청 공공 API에서 수집한 설화·민담 505건을 ChromaDB에
1{,}749 청크로 벡터화하여 적재하되, 비짓제주 관광지와 매핑된 397 설화·민담만
동행자 채팅 검색 대상으로 한정하였다. 한국관광공사 KTO 분류와 키워드 규칙을 결합한
3단계 필터로 Visit Jeju 장소 3{,}284개에서 관광지 820개를 선별하였고, 그중 524개
관광지가 397 설화·민담과 직접 매핑되어 사전 매핑 테이블 6{,}726건을 구성한다(한
설화가 평균 17개 장소에 등장). 장소-설화 사전 매핑 30{,}947건 중 GPS 보조 매핑을
제외한 정합성 매핑 6{,}726건만 코스 추천에 사용하여, 두 검색 경로 모두에 일관된
\emph{장소 정합성 원칙}을 적용하였다(원본 매핑 134k 대비 노이즈 $-77\%$). 최종 보고 시점 기준으로 홈 화면,
2단계 코스 추천, GPS 탐험, 5종 페르소나 동행자·TTS, 장소 리뷰, 설화 3단계
가공 뷰어, 일지·민화 자동 생성, 세션 복원, ``내 코스'' 자동 이동 등 시연
플로우 10단계 전 구간이 구현 완료되었다.}

\vspace{6pt}
{\small\noindent\textbf{키워드}\quad
RAG, 페르소나 LLM, ChromaDB, LangGraph, iOS Swift, 위치 기반 서비스,
제주 설화·민담, 멀티모달 생성, TTS, 여행 동행자}

\vspace{8pt}

% ── 1. 서론 ─────────────────────────────────────────────
\section{서~~~론}

제주도는 연간 1{,}500만 명 이상이 방문하는 국내 최대 관광지이나, 현행 관광
앱의 대부분은 맛집·숙박·교통 중심의 정보 제공에 머물고 있다. 제주도에는
설화 182건, 민담 323건 등 505건의 문화유산 콘텐츠가 공공 데이터로 개방되어
있음에도, 이를 관광 경험과 결합한 서비스는 부재한 상황이다.

최근 LLM(Large Language Model)의 발전과 함께, 특정 도메인 지식을 벡터
데이터베이스에 저장하고 검색·생성을 결합하는 RAG 기술이 주목받고 있다.
RAG는 모델의 환각(Hallucination) 문제를 억제하면서 도메인 특화 지식을
활용할 수 있어, 정확성이 요구되는 문화유산 콘텐츠 분야에 적합하다. 또한
LLM 워크플로우 프레임워크와 멀티모달 생성 모델의 발전은 단순한 정보
조회를 넘어 ``계획-탐험-기록''의 전 여정에 AI가 개입할 수 있는 기술적
토대를 제공한다.

본 연구는 \textbf{``실제 검증된 여행 코스 위에 설화를 얹는다''} 는 컨셉으로,
제주관광공사 Visit Jeju의 실제 관광객 여행 일정 데이터를 기반으로 LangGraph
스타일 LLM 함수 파이프라인이 설화·민담이 풍부한 코스를 추천하고, 여행
중에는 5종의 페르소나 기반 AI 설화 동행자가 장소에 깃든 이야기를 텍스트와
음성으로 전달하며, 여행 후에는 일지와 민화 이미지가 자동 생성되는 iOS 앱을
구현하는 것을 목표로 한다.

본 연구의 졸업프로젝트로서의 학문적·실무적 기여는 다음 세 가지로 정리된다.
첫째, \emph{Visit Jeju 실제 관광 코스 데이터와 KTO 분류를 결합}한 신뢰성
있는 코스 큐레이션 위에 무형 문화유산을 레이어로 결합하는 역방향 매핑
구조를 제안하여, LLM 자유 생성 코스의 비현실성과 단순 정보 나열식 가이드
앱의 빈약함을 동시에 해결하였다. 둘째, \emph{5종 페르소나 LLM 챗봇과 RAG의
결합}을 통해 같은 설화를 다양한 캐릭터로 재해석하여, 일반 챗봇이 결여한
캐릭터 일관성과 현지감을 구현하였다. 셋째, \emph{GPT-4o 단일 호출의 구조화
출력으로 일지와 민화 이미지 프롬프트를 동시 생성}하는 멀티모달 파이프라인을
제시하여 사용자별 고유 결과물을 합리적 지연·비용 내에서 제공하였다.

% ── 2. 관련 연구 ───────────────────────────────────────
\section{관련 연구}

\subsection{RAG(Retrieval-Augmented Generation)}

RAG는 Lewis et al.(2020)이 제안한 방법으로\,[1], 외부 지식 저장소에서 관련
문서를 검색(Retrieve)한 뒤 이를 문맥으로 제공하여 LLM이 더 정확한 답변을
생성(Generate)하도록 한다. 벡터 임베딩 기반의 의미론적 검색을 통해 키워드
검색의 한계를 극복하며, 특정 도메인에 특화된 지식베이스 구축에 활용된다.
본 연구는 RAG를 활용하여 ChromaDB\,[5]에 저장된 제주 설화·민담을 실시간으로
검색하고, 여행 현장 컨텍스트와 페르소나 말투에 맞는 해설을 생성하는 데
적용한다.

\subsection{위치 기반 문화 콘텐츠 서비스}

Pok\'{e}mon GO, Ingress 등의 위치 기반 서비스는 GPS를 활용하여 사용자가
실제 장소를 방문하도록 유도한 바 있다. 문화유산 분야에서는 Google Arts \&
Culture, 국내 문화재청의 e뮤지엄 등이 위치 기반 콘텐츠를 제공하나, 지역
설화·민담과 같은 무형 문화유산과 AI 코스 추천을 결합한 서비스는 미흡한
실정이다. 본 연구는 실제 관광객 동선 데이터와 무형 문화유산을 결합하여,
기존 서비스에서 미흡했던 설화·민담 기반 AI 코스 추천을 구현한다.

\subsection{대화형 AI 인터페이스}

ChatGPT, Claude 등의 등장 이후 사용자가 자연어로 정보를 탐색하는 대화형
인터페이스가 보편화되고 있다. RAG를 결합한 도메인 특화 챗봇은 일반 LLM의
지식 한계를 보완하면서 특정 주제에 깊이 있는 답변을 제공한다. 본 연구에서는
이를 관광 맥락에 적용하여 현장형 AI 설화 동행자를 구현한다.

\subsection{LLM 에이전트 및 LangGraph}

Yao et al.(2023)이 제안한 ReAct(Reasoning + Acting) 패턴은\,[8] LLM이 추론
단계와 도구 사용 단계를 반복하며 복잡한 작업을 수행하도록 한다. LangGraph\,[9]는
이러한 LLM 워크플로우를 유향 그래프(StateGraph)로 명시적으로 정의하여,
조건부 분기·루프·병렬 실행을 지원한다. 본 연구에서는 LangGraph의
\texttt{invoke(state) → state} 인터페이스를 모방한 함수 파이프라인을 채택하여
2단계 코스 추천을 구조화하였으며, ReAct 패턴은 사용하지 않는다. 코스 추천·
동행자 채팅·일지 생성 등 응답 지연이 사용자 이탈에 직결되는 핵심 경험에는
결정론적 함수와 단일 LLM 호출의 조합이 도구 자율 호출 라운드를 누적시키는
ReAct보다 비용·지연·환각 통제 측면에서 우월하다고 판단하였다.

\subsection{페르소나 기반 대화 인터페이스}

Park et al.(2023)은 ``Generative Agents''에서 LLM이 기억·반영·계획 모듈을
통해 일관된 개성을 유지하며 상호작용하는 시뮬레이션을 제시하였다\,[11].
이후 캐릭터 일관성·말투·세계관을 시스템 프롬프트에 인코딩하여 유지하는
페르소나 챗봇 연구가 활발히 진행되고 있다. 일반 LLM 챗봇은 ``유용한 비서''
어조의 단일 페르소나에 수렴하기 쉬워, 문화 콘텐츠 전달에서 요구되는
캐릭터별 정체성·방언·세계관 차이를 충분히 표현하지 못한다. 본 연구는
설화 유형별로 5종의 페르소나를 분리하여 \emph{독립된 말투 가이드와 퓨샷
예시}를 시스템 프롬프트에 인코딩하고, 사용자의 취향 카테고리 점수에 따라
자동 배정함으로써 페르소나의 일관성과 적합성을 동시에 확보한다.

\subsection{멀티모달 콘텐츠 생성 및 구조화 출력}

GPT-4o와 같은 멀티모달 LLM은 텍스트뿐 아니라 이미지·오디오를 통합적으로
다룬다\,[12]. 또한 함수 호출(function calling) 및 JSON 스키마 응답
기능은 LLM의 출력을 다운스트림 시스템이 안정적으로 소비할 수 있는 구조화된
형태로 강제한다. 이미지 생성 모델인 gpt-image-1\,[13]은 텍스트 프롬프트로
일관된 화풍의 이미지를 생성하며, 텍스트 일지와 이미지 프롬프트를 단일 호출의
구조화 출력으로 동시에 받아 다음 단계에 그대로 전달할 수 있다. 본 연구는
이 패턴을 활용하여 LLM 호출 횟수와 지연을 최소화하면서 일지와 민화 이미지를
함께 생성한다.

% ── 3. 시스템 설계 ────────────────────────────────────
\section{시스템 설계}

\subsection{전체 아키텍처}

본 시스템은 iOS 앱(Swift/SwiftUI), FastAPI 백엔드, 데이터 레이어의 3계층에
Firebase 인증·동기화 레이어가 결합된 구조로 구성된다. iOS 앱은 REST/SSE를 통해
백엔드와 통신하며, 백엔드는 2단계 코스 추천 파이프라인,
설화 동행자 SSE 엔드포인트, GPS 핀 조회 API, TTS API,
장소 리뷰 API, 일지·이미지 생성 API, 설화 3단계 가공 캐시 API로 구성되며, 사용자
식별이 필요한 엔드포인트는 Firebase ID Token을 검증한다.
데이터 레이어는 Visit Jeju 여행일정 DB(큐레이션 코스)\,[4], ChromaDB(설화·민담
벡터 DB)\,[5], SQLite(메타데이터·GPS 좌표·리뷰·설화 가공 캐시)로 이루어지며, 사용자
계정·로그인 세션과 저장 코스의 다기기 동기화는 Firebase Authentication(Google
Sign-In)과 Firestore가 담당한다.

\begin{footnotesize}
\begin{verbatim}
┌──────────── iOS App (Swift/SwiftUI) ────────────┐
│ Home │ Course Wizard │ Explore │ MyCourse │Sheet│
└──────┬─────────────┬──────────┬───────────┬─────┘
       │REST         │REST      │SSE        │REST
       ▼             ▼          ▼           ▼
┌───────────────── FastAPI Backend ───────────────┐
│ /course/list,detail   /travel/companion (SSE)   │
│ /pins  /tts  /place/review  /tourist/info       │
│ /folklore/connection  /folklore/story           │
│ /travel/journal (+image)                        │
└──┬──────────┬──────────┬──────────┬─────────────┘
   │          │          │          │
   ▼          ▼          ▼          ▼
ChromaDB     SQLite     OpenAI    Google     Firebase
(1,749 적재  courses    GPT-4o /  Maps SDK   Auth +
 /매핑 청크  9,134 +    gpt-      / KTO      Firestore
 검색)       place_rev  image-1   TourAPI    (iOS SDK)
 chunks
\end{verbatim}
\end{footnotesize}

\begin{center}
{\small 그림 1~~전체 시스템 아키텍처}
\end{center}

사용자 흐름은 \emph{계획·탐험·기록}의 단일 여정으로 구성된다. 홈에서는
오늘의 설화 카드, 제주 지도와 설화 핀, 추천 코스 3선이 첫 화면에 노출되어
앱의 정체성을 즉시 전달한다. 계획 단계에서는 취향(지역·스타일·기간·설화취향)을
입력하면 1순위 코스가 즉시 상세 미리보기로 표시되며, ``새로운 추천'' 버튼으로
2·3순위 코스를 순서대로 확인할 수 있다. 탐험 단계에서는 실시간 GPS 안내와
함께 장소 도착 시 AI 설화 동행자가 말을 건다. 기록 단계에서는 여행 종료
후 GPT-4o가 일지와 민화 이미지를 자동 생성한다. 앱 재시작 시 SessionRestoreView가
진행 중인 탐험 세션을 자동 복원하여 앱 강제 종료에도 여행 연속성을 보장한다.

\subsection{데이터 파이프라인}

\textbf{Visit Jeju 여행 데이터.}\quad 제주관광공사 Visit Jeju 공공데이터\,[4]에서
여행장소(GPS 좌표 포함)와 여행세부일정을 수집하였다. 장소명 기준 조인으로
일정별 방문 장소의 GPS 좌표를 확보한 뒤, 공항·항구 중복 도배 코스 등
저품질 코스를 필터링하여 큐레이션 코스 DB(9{,}134개 코스, 146{,}508개
코스-장소 레코드)를 구성하였다.

\textbf{비짓제주 장소 분류 --- KTO + 키워드 + 후처리 3단계 필터링.}\quad
Visit Jeju 원본 데이터에는 장소 종류(관광지·식당·숙박 등) 정보가 없어,
초기에는 모든 장소가 설화 매핑 대상이 되었다. 이로 인해 약 134{,}000건의
장소-설화 매핑 중 상당수가 ``식당 ↔ 설화'' 같은 무의미한 매핑으로
구성되었다. 이를 해결하기 위해 다음 3단계 필터를 도입하였다.

\begin{enumerate}
\setlength{\itemsep}{0pt}
\item \textbf{KTO 정식 관광지 매칭}: 한국관광공사 TourAPI\,[10]의 제주
공식 관광지 약 618개에 대해 GPS 거리·이름 유사도로 비짓제주 장소와
매칭하여 약 456개를 라벨링하였다.
\item \textbf{키워드 기반 분류}: 장소명에 ``오름·해변·박물관·4.3·사격장''
등 지정 패턴이 포함되면 자연·문화·유적·체험 카테고리로 자동 분류하여
약 1{,}400개를 추가 라벨링하였다.
\item \textbf{후처리 강제 재분류}: ``○○점·갈비·회국수'' 등 식당
패턴이 관광지로 잘못 분류된 경우를 보정하였다(29건).
\end{enumerate}

결과적으로 3{,}284개 장소를 12개 카테고리로 정리하고, 관광지(자연/문화/유적/
체험/마을시설/종교) 820개만 설화 매핑 대상으로 한정하였다. 그중 524개 관광지가
설화·민담과 실제로 매핑되어 사전 매핑 테이블 30{,}947건을 구성한다. 표 1은
필터 전후의 매핑 신뢰성 변화를 보여준다.

\begin{table}[H]
\centering
{\small 표 1~~장소 분류 필터의 매핑 신뢰성 효과}
\label{tab:filter}
\vspace{2pt}

\small
\begin{tabular}{lrr}
\toprule
지표 & 분류 전 & 분류 후 \\
\midrule
매핑 대상 장소 & 3{,}284 & \textbf{820} \\
매핑 보유 관광지 & --- & \textbf{524} \\
장소-설화 매핑 건수 & $\sim$134{,}000 & \textbf{30{,}947} \\
GPS 보조 매핑 노이즈 & $\sim$112{,}000 & $\sim$24{,}221 \\
신뢰성 개선(노이즈 제거) & --- & \textbf{$-77\%$} \\
\bottomrule
\end{tabular}
\end{table}

\textbf{설화·민담 → 비짓제주 관광지 매핑.}\quad 설화\,[2]는 텍스트 내 지명을 제주
행정구역 화이트리스트 기반 매칭으로 추출하고, 민담\,[3]은 채록 메타데이터의 조사
장소를 파싱하여 비짓제주 관광지(이름·주소)와 직접 매칭한다.
설화·민담 505건 전체를 OpenAI Text Embedding\,[6]으로 청크 분할하여 ChromaDB에
1{,}749 청크로 벡터화하였고, 그중 비짓제주 관광지와 직접 매핑된 항목(505건 중 397건,
약 78\%)은 \texttt{is\_mapped=True} 메타데이터를 부여하여 동행자 채팅의 의미 검색
대상이 된다. 매핑이 발생하지 않은 108건(본문이 매우 짧거나 매칭 가능 지명이 없는
경우)도 ChromaDB에 함께 적재되어 향후 자유 발화 챗봇 등 다른 기능에서 재활용 가능하도록
보존하였다.

\subsection{홈 화면 설계}

홈 화면은 앱 진입 시 ``이 앱이 무엇을 하는 서비스인가''를 5초 안에 전달하는
역할을 한다. 초기 설계에서는 ``코스 만들기''와 ``내 코스''만 탭으로 존재했으나,
두 탭 모두 사용자가 명확한 의도를 가지고 누르는 액션 탭이지 메인 화면이
아니었다. 이에 ScrollView 기반으로 다음 3개 요소를 수직 배치한 홈 화면을
신설하였다.

\begin{enumerate}
\setlength{\itemsep}{0pt}
\item \textbf{오늘의 설화 카드}: 현재 날짜를 시드로 해시하여 매일 다른 설화를
일관되게 노출. 동일 일자 내에는 항상 동일한 설화가 표시되어 ``오늘의 설화''
정체성을 유지한다.
\item \textbf{제주 지도 + 설화 핀}: 380pt 고정 높이의 지도 위에 설화
위치 핀을 표시하여 무형 문화유산의 공간 분포를 직관적으로 전달한다.
\item \textbf{추천 코스 3선}: 별도 취향 입력 없이도 대표 코스 3개를
가로 카드로 노출하여 ``코스 만들기'' 위저드의 진입 부담을 낮춘다.
\end{enumerate}

\subsection{코스 추천: 2단계 함수 파이프라인과 가중 점수 기반 후보 풀}

응답 지연을 최소화하기 위해 코스 추천을 두 단계의 함수 파이프라인으로
분리하였다. 사용자 입력은 \emph{region}(북부/남부/동부/서부/\textbf{전체} 5지선다),
\emph{category\_scores}(메인 카테고리 3점·사이드 카테고리 1점), \emph{duration\_days}의
3가지이다.

\textbf{1단계 후보 풀 구성과 가중 점수 계산(\texttt{POST /course/list}).}\quad
\texttt{course\_list\_agent} 모듈이 다음 결정론적 흐름을 \emph{LLM 호출 0회}로
수행한다.

\begin{enumerate}
\setlength{\itemsep}{0pt}
\item SQLite \texttt{curated\_courses}(전체 1{,}090개 풀)에서 region·duration
조건으로 후보 50개를 추출한다. region이 ``전체''이면 region 필터를 풀어
1{,}090개 모든 코스가 후보 풀에 포함되고, 4지선다(북부/남부/동부/서부) 선택
시에는 요청 지역 + ``전체'' 태그 코스(약 297개)에서 후보를 잡는다.
\item 각 후보 코스에 대해 \texttt{course\_places JOIN place\_folklore\_mapping}을
수행하되, \textbf{specificity $\ge 5$}이고 \textbf{\texttt{source != 'gps\_assist'}}인
매핑만 사용한다(아래 정합성 원칙 참조).
\item specificity 가중치를 $w = \mathrm{specificity}/5.0$으로 정의하고, 카테고리
점유율 $\mathrm{share}_c = w_c / \sum w$를 구한 뒤,
$\mathrm{score} = \big(\sum_c \mathrm{category\_scores}[c] \times \mathrm{share}_c\big) \times 10$
으로 코스 점수를 산정한다.
\item $\mathrm{score} > 0$인 코스만 풀에 남기고 \texttt{random.sample(3)}으로
\textbf{매번 무작위 3건}을 뽑는다. 시연·실사용 모두에서 동일 입력에 동일 결과가
나오는 경직성을 해소한다.
\end{enumerate}

\textbf{2단계 상세화(\texttt{POST /course/detail}).}\quad 결정론적 함수가
선택된 코스의 모든 장소에 대해 사전 매핑 테이블에서 설화·민담 후보를 가져온다
(이때도 동일하게 \texttt{source != 'gps\_assist'} 필터 적용). 시연 응답 속도와
비용 절감을 위해 본 단계의 LLM 내러티브 생성은 \emph{비활성화}되어 있으며,
필요 시 \texttt{with\_structured\_output} 스키마로 2문장 여행 내러티브를 생성하는
경로를 함수 토글로 보존하였다.

\textbf{랜덤 풀의 의미.}\quad 카테고리 매칭이 전혀 안 되는 코스는 풀에서
배제하므로 추천의 의미는 유지하면서, 같은 취향으로 다시 진입해도 매번 다른
조합을 본다. score $> 0$ 코스가 0건이면 빈 결과를 반환하여 클라이언트가
``조건에 맞는 코스 없음'' 화면으로 자연스럽게 유도된다(점수 무관 fallback 제거).

\subsection{장소 정합성 원칙 --- 두 검색 경로의 일관된 패턴}

본 시스템은 \emph{``장소와 진짜 관련 있는 설화만 사용자에게 노출한다''}는
원칙을 의미 검색(ChromaDB)과 사전 매핑(SQLite)이라는 \textbf{두 검색 경로에
일관되게 적용}한다. 한 학기 운영 중 발견한 ``근처에 있다 $\ne$ 그 명소의
설화''라는 정합성 우려를 동일한 설계 패턴으로 해소한 결과이다.

\begin{table}[H]
\centering
{\small 표 7~~두 검색 경로의 정합성 필터}
\label{tab:integrity}
\vspace{2pt}

\footnotesize
\begin{tabularx}{\columnwidth}{p{0.27\columnwidth}X}
\toprule
경로 & 메타데이터/컬럼 \& 쿼리 필터 \\
\midrule
ChromaDB 의미 검색
(동행자 채팅) &
메타데이터 \texttt{is\_mapped} (bool) /
쿼리 \texttt{where=\{is\_mapped: True\}} \\
\addlinespace[2pt]
사전 매핑 테이블
(코스 추천) &
컬럼 \texttt{source} $\in$ \{\texttt{name\_match},
\texttt{collection\_site}, \texttt{gps\_assist}\} /
쿼리 \texttt{WHERE source != 'gps\_assist'} \\
\bottomrule
\end{tabularx}
\end{table}

(1) \textbf{ChromaDB 의미 검색}: 모든 설화·민담 505건 1{,}749 청크를 적재하되,
비짓제주 관광지와 직접 매핑된 397건에만 \texttt{is\_mapped=True} 메타데이터를 부여하고
동행자 채팅 검색 쿼리에 \texttt{where} 필터를 적용해 \emph{매핑된 설화 청크만} 검색
대상으로 제한한다. (2) \textbf{사전 매핑 테이블}: 매핑 출처를 \texttt{source} 컬럼
(\texttt{name\_match}, \texttt{collection\_site}, \texttt{gps\_assist})으로 명시한 뒤,
코스 추천 쿼리에 \texttt{source != 'gps\_assist'} 필터를 적용해 8\,km 반경 GPS 보조
매핑(24{,}221건)을 제외하고 지명 텍스트 매칭과 민담 채록 장소 매칭 \textbf{6{,}726건}만
사용자에게 노출한다.

두 경로 모두 데이터 자체는 그대로 보존되므로, 향후 정책 변경(예: 보조 매핑
일부 허용) 시 코드 한 줄의 필터 토글로 대응 가능하다.

\subsection{장소 상세 --- 설화 데이터 3단계 가공}

장소 상세 화면의 설화는 원본이 학술적 어투의 채록 본풀이 텍스트로,
\texttt{C\_M\_022}와 같은 내부 코드가 노출되고 긴 본문이 한꺼번에 펼쳐져
여행 중 5초 안에 호기심을 끌기 어려웠다. 이를 다음 3단계로 점진적으로
가공하였다.

\begin{table}[H]
\centering
{\small 표 2~~설화 데이터 3단계 가공}
\label{tab:lv}
\vspace{2pt}

\footnotesize
\begin{tabularx}{\columnwidth}{lX}
\toprule
단계 & 내용 \\
\midrule
Lv1 & 카드에서 코드 접두사 제거. LLM이 생성한 호기심 자극형 후크 1줄을
노출(예: ``심방으로 살다 다시 이승으로 돌아온 굴묵밭 할망''). 원본 요약은
접어두기. \\
Lv2 & 카드 상단에 ``이 마을의 본향당 신을 모시는 이야기''와 같이 \emph{장소와
설화의 관계 1줄}을 추가. 백엔드 LLM이 장소명+설화 본문을 받아 생성. \\
Lv3 & 카드 탭 시 인스타그램 스토리 느낌의 풀스크린 뷰어 진입. 5--7페이지의
동화체 리라이팅 스토리, 좌우 스와이프, 페이지 인디케이터. \\
\bottomrule
\end{tabularx}
\end{table}

\textbf{영구 캐시 테이블.}\quad Lv2·Lv3는 LLM 호출이 필요하지만, 동일 (장소,
설화) 조합은 항상 동일한 결과를 갖는다. 이에 백엔드에 \texttt{folklore\_connection\_cache}
(Lv2 연결 문장)와 \texttt{folklore\_story\_cache}(Lv3 동화체 스토리) 두 테이블을
신설하여 \emph{한 번 생성된 결과를 영구 저장}한다. 동일 키 재호출 시 LLM을
호출하지 않고 캐시에서 즉시 반환하여 비용과 지연을 모두 최소화한다.

\subsection{AI 설화 동행자: 5종 페르소나 시스템}

여행 중 경험의 중심축으로 5종의 페르소나 기반 AI 설화 동행자를 구현하였다.
본 시스템의 핵심 경험인 AI 설화 동행자는 \emph{단일 LLM 호출에 페르소나
시스템 프롬프트와 ChromaDB RAG 컨텍스트를 결합}한 구조이며, 도구 자율 호출
ReAct 패턴은 사용하지 않는다. 페르소나 정의(이름·말투·어휘·퓨샷 예시)는
별도 설정으로 관리하고, 동행자는 해당 장소 주변 설화를 ChromaDB에서 벡터
검색하여 맥락으로 주입받으며, FastAPI SSE를 통해 토큰 단위 스트리밍으로
응답한다.

\begin{table}[H]
\centering
{\small 표 3~~AI 설화 동행자 5종 페르소나}
\label{tab:persona}
\vspace{2pt}

\footnotesize
\begin{tabularx}{\columnwidth}{lXl}
\toprule
페르소나 & 설화 유형 & 말투 핵심 \\
\midrule
마을 할망       & 마을 공동체 전승       & 따뜻한 제주 사투리 \\
당신·심방       & 무속신화·신격 전승     & 고어체 신탁 \\
영등신·해녀     & 해양·어촌 전승         & 직선적 반말, 해녀 어휘 \\
도깨비         & 생활민담·교훈담         & 장난기, 질문 되받기 \\
도체비         & 초자연 존재담           & 문장 끊기, 자기 부정 \\
\bottomrule
\end{tabularx}
\end{table}

\textbf{독립 프롬프팅 --- 공유 사투리 가이드 배제.}\quad 캐릭터마다 말투
가이드와 4--5개의 퓨샷 예시를 \emph{독립적으로} 설계하였다. 동일한 제주
사투리 가이드를 공유하지 않고, 각 캐릭터의 설화 유형에 맞는 전용 어휘와
발화 패턴을 별도로 정의한다. 예를 들어 해녀 페르소나는 ``물질·테왁·숨비소리·
바릇·불턱·영등'' 등 6개의 해녀 전용 어휘 사전과 ``시원시원하고 직선적,
돌려 말하지 않음''의 규칙을 갖는 반면, 심방 페르소나는 ``\textasciitilde
이니라/하시라'' 고어체 어미를 사용하며 사투리는 최소화한다.

\textbf{페르소나 자동 배정 + 수동 변경.}\quad 사용자의 취향 카테고리 최고
점수가 ``무속신화''면 심방, ``해양·어촌''이면 해녀 등 자동 배정 규칙으로
초기 동행자가 결정된다. 코스 미리보기 화면에 배정된 동행자 이름과
``바꾸기'' 버튼이 표시되며, 탭하면 \texttt{CompanionPickerSheet} 바텀시트로
5개 카드 중 직접 선택할 수 있다. 변경 결과는 탐험 진행 화면에 override로
전달된다.

\subsection{페르소나 TTS 시스템}

5종 페르소나의 캐릭터성은 텍스트만으로 절반밖에 살지 않는다. 이에 페르소나
전용 TTS 음성을 도입하여 텍스트와 음성을 동시에 출력한다.

\begin{table}[H]
\centering
{\small 표 4~~페르소나별 TTS 음성 색깔}
\label{tab:tts-voice}
\vspace{2pt}

\footnotesize
\begin{tabularx}{\columnwidth}{lX}
\toprule
페르소나 & 음성 색깔 \\
\midrule
당신·심방 & 저음 남자, 위엄 있고 신비함 \\
도깨비   & 스토리텔러 톤, 경쾌·익살 \\
마을 할망 & 부드러운 여자, 정겨움 \\
영등·해녀 & 강한 여자, 시원한 바다 톤 \\
도체비   & 차분한 남자, 속삭이듯 신비함 \\
\bottomrule
\end{tabularx}
\end{table}

\textbf{TTS 재생 3 포인트.}\quad (1) 장소 도착 오버레이가 뜨자마자 페르소나
자기 목소리로 자동 인사, (2) 채팅에서 AI 응답 시 글자가 한 자씩 노출되는
동시에 페르소나 목소리로 발화, (3) 여행 일지 화면에서 재생 버튼 탭 시
일지 전체를 차분한 내레이터 톤으로 낭독.

\textbf{글자/음성 타이밍 동기화.}\quad 핵심 규칙은 ``글자 출력과 음성을
거의 같은 타이밍에 끝나게 한다''이다. 글자 다 보인 뒤 1초 멍하니 기다리고
음성이 시작되는 어색함을 제거하기 위해, 글자 노출 속도를 TTS 길이에 맞춰
조정한다. 우측 상단 스피커 아이콘 한 번 탭으로 전체 음소거가 가능하며,
설정은 앱 재실행 후에도 유지된다.

\subsection{GPS 탐험 모드와 Google Maps SDK 전환}

iOS CoreLocation 기반 지오펜스로 도착을 감지한다. 도보 모드에서는 반경
100m 진입 후 30초 체류, 차량 모드에서는 반경 300m 진입 후 30초 체류 시
도착으로 판정한다. 30초 체류 조건은 잠깐 지나치는 오탐을 방지한다. 도착
판정 시 동행자 오버레이가 전체화면으로 등장하며, 사용자가 탭하면 해당
장소의 설화를 주제로 한 채팅이 시작된다.

\textbf{Apple Maps → Google Maps SDK 전환.}\quad 초기에는 Apple Maps를
사용하였으나, (1) 임베드 미리보기에서 Apple 워터마크가 시연 시 그대로
노출되고, (2) 길찾기가 Apple Maps 앱으로 빠지며, (3) 한국 사용자 다수가
구글 지도를 쓰는 실사용 환경과 어긋난다는 문제가 있었다. GoogleMaps iOS
SDK\,[7]를 도입하여 임베드 미리보기와 길찾기 양쪽을 Google Maps로 통일하였고,
구글맵 앱 미설치 시 \texttt{https://www.google.com/maps/dir/...} 웹 URL로
자동 폴백한다. 장소 상세 정보는 \texttt{GET /tourist/info}를 통해 KTO
TourAPI\,[10]에서 사진·정보를 7일 캐시 기준으로 조회한다.

\subsection{장소 리뷰 시스템}

장소에 대한 다른 여행자들의 반응을 보여주는 커뮤니티 섹션을 구현하였다.
설화 상세 화면 하단에 5개 태그(소름 돋아요·감동이에요·신기해요·무서워요·
역사적이에요)의 통계 막대와 최근 감상 노트 3개가 노출된다. 데이터가 없으면
섹션 자체를 숨겨 빈 상태를 노출하지 않는다.

\textbf{자동 바텀시트.}\quad 탐험 중 GPS 도착 → 동행자 채팅 → 채팅 닫기의
흐름이 끝나면 \texttt{PlaceReviewSheet}가 자동으로 올라온다. 태그 1개 이상
선택 시 ``남기기'' 버튼이 활성화되고, ``건너뛰기''로 닫아도 된다. 리뷰는
Firebase Authentication 로그인 사용자의 경우 \texttt{user\_id}(Firebase uid)로,
미로그인 사용자의 경우 익명 \texttt{device\_id}(UUID)로 식별되어 저장되며,
같은 사용자가 동일 장소를 재방문하면 덮어쓴다. 제출 실패는 조용히 처리되어 앱 동작에 영향을 주지 않는다.

\textbf{일지 생성 시 GPT 프롬프트 주입.}\quad 제출된 리뷰는 \emph{여행 일지
생성 시 GPT 프롬프트에 자동 주입}되어, 커뮤니티 반응이 일지 이야기 흐름에
자연스럽게 녹는다. 이로써 같은 장소를 방문해도 사용자별 리뷰 선택에 따라
일지 분위기가 달라진다.

\subsection{여행 일지·민화 이미지 동시 생성}

탐험 종료 후 방문 장소 요약 화면을 거쳐 여행 일지를 자동 생성한다. 입력은
\emph{방문 장소 시퀀스·동행자 대화 로그·커뮤니티 반응 통계} 3가지이며,
GPT-4o의 \textbf{단일 호출에서 구조화 출력(JSON 스키마)}으로 일지 텍스트와
민화 이미지 프롬프트를 동시에 받는다. 이미지 프롬프트는 gpt-image-1\,[13]에
그대로 전달되어 한 장의 민화 이미지로 렌더링된다.

\textbf{이미지 프롬프트 규칙(요약).}\quad (1) 일지의 가장 인상 깊은 한 장면을
영어 80--150단어로 묘사, (2) ``Korean traditional minhwa style, flat color,
thick black outline, decorative composition, paper-textured background...''로
반드시 시작, (3) 장소(오름·해안·당·해녀 잠수)·인물(심방 할머니·해녀·도깨비
영혼)·분위기·상징(소나무·매화·파도·구름·산) 요소를 포함, (4) 사람 얼굴
디테일 묘사 금지(민화의 단순화 표현 유지), (5) 영문 내 한글/한자 글씨 금지,
(6) 폭력적 표현 금지.

\textbf{단일 호출 동시 생성의 효과.}\quad 일지와 이미지 프롬프트를 별도
호출로 처리할 경우 LLM 호출 2회와 그에 따른 비용·지연이 발생하지만, JSON
스키마 강제로 단일 호출에서 동시에 받으면 호출 1회로 완결된다. 이미지
생성이 실패해도 일지 텍스트는 정상 표시되어 부분 실패에도 견고하다. 전체
응답 시간은 약 25--30초이며, 생성된 일지는 iOS ShareLink를 통해 외부
공유가 가능하다.

\subsection{세션 복원과 여행 완료 동선}

\textbf{세션 복원.}\quad 탐험 도중 앱이 강제 종료되거나 디바이스가 재부팅되어도
다음 실행 시 SessionRestoreView가 자동으로 진행 중 세션을 감지하여
사용자에게 ``이전 여행을 계속하시겠습니까?'' 프롬프트를 표시한다. 세션 ID·
현재 day·다음 방문 장소·동행자 상태·대화 히스토리를 모두 복원하여 여행
연속성을 보장한다.

\textbf{여행 완료 동선.}\quad 일지 ``완료'' 탭 시 즉시 ``내 코스'' 탭으로
이동시키고, 그 안에서 방금 끝낸 코스를 자동 스크롤하여 1.5초간 배경 강조로
하이라이트한다. 이전 구현에서는 코스 만들기 도중 노출되던 추천 리스트
화면으로 빠져 마무리 임팩트가 사라지는 문제가 있었다. 본 동선은 \emph{일지
완료 직후 앱이 종료되어도} 다음 실행 시 ``내 코스'' 탭이 자동 선택되고
하이라이트가 보존되도록 견고하게 구현되었다.

% ── 4. 설계 변경 이력 ──────────────────────────────────
\section{설계 변경 이력}

\subsection{플랫폼: Flutter → iOS Swift/SwiftUI}

기초 조사서에는 Flutter 크로스플랫폼으로 계획하였으나, 이전 산책 앱 개발
경험에서 Flutter의 위치 기반 기능 일부를 결국 Kotlin·Swift 네이티브로
별도 구현해야 했던 교훈을 반영하여 iOS 네이티브 개발로 전환하였다.
CoreLocation 지오펜스, 백그라운드 GPS, iOS 알림 딥링크 등 위치 기반 기능의
정확도와 통합성이 향상되었다.

\subsection{코스 추천 방식: 3차 재설계}

코스 추천 로직은 세 차례의 재설계를 거쳤다. \emph{1차(설화→코스 역설계)}는
설화 6노드 LangGraph로 설화 배경 장소를 먼저 추출하고 Visit Jeju로 동선을
보완하는 방식이었으나, 설화 위치 분포가 희박할 때 코스가 비현실적으로
구성되는 문제가 있었다. \emph{2차(설화 취향 점수)}는 사용자 설화 취향을
내부 7개 카테고리 점수로 변환하여 필터링하는 구조였으나, 카테고리 분류가
실제 설화 데이터 다양성을 반영하지 못하였다. \emph{3차(코스→설화 매핑,
현재)}는 사용자가 먼저 정하는 지역·기간을 입력받아 Visit Jeju 실제 코스를
먼저 선별하고 그 위에 설화를 레이어로 얹는 방식으로, 코스의 현실성은 실제
데이터가 보장하고 설화는 코스를 풍부하게 만드는 역할로 분리되었다.

\subsection{지도 SDK: Apple Maps → Google Maps}

장소 상세의 임베드 미리보기와 길찾기 모두 Google Maps iOS SDK로 통일하였다.
한국 사용자 다수의 실사용 환경에 부합하며, 시연 시 Apple 워터마크 노출
문제도 해소되었다. 미설치 환경 대비 웹 URL 폴백을 함께 구현하였다.

\subsection{사용자 입력 UI: 16조합 변환표 → 5개 카테고리 직접 선택}

기존에는 ``이야기 결 4지선다 + 배경 4지선다''로 16조합 입력을 받아 내부
변환표로 5개 설화 카테고리 점수로 환산하였다. 발표 시 변환표 그림까지
설명해야 하는 부담이 있었고, 사용자 입장에서도 자신의 입력이 어떻게
해석되는지 불투명하였다. 이에 5개 설화 카테고리를 친근한 텍스트로 직접
노출하여 \emph{메인 1개 + 사이드 1개}를 선택하게 하고 점수는 메인 3점,
사이드 1점으로 단순화하였다. 발표 슬라이드 한 줄로 설명이 끝난다.

\subsection{장소 매핑: 전체 → KTO 관광지 분류 820개}

3.2절에서 상술한 KTO + 키워드 + 후처리 3단계 필터를 도입하여, 3{,}284개
장소 중 관광지 820개만 설화 매핑 대상으로 한정하였다. 식당·카페·숙박·
상점에 설화 핀이 붙던 신뢰성 문제가 해소되었다.

% ── 5. 구현 결과 ─────────────────────────────────────────
\section{구현 결과}

최종 보고 시점 기준으로 시연 플로우 10단계 전 구간이 구현 완료되었다.
주요 기능은 표 5에 정리하였다.

\begin{table}[H]
\centering
{\small 표 5~~구현 완료 기능 요약}
\label{tab:impl}
\vspace{2pt}

\footnotesize
\begin{tabularx}{\columnwidth}{lX}
\toprule
영역 & 구현 내용 \\
\midrule
홈 & 오늘의 설화(날짜 시드) + 지도/핀(380pt) + 추천 코스 3선 \\
취향 입력 & 지역·기간·5카테고리(메인3+사이드1) 위저드 \\
코스 추천 & 2단계 API + 후보 풀 랜덤 3건, 빈 결과 자연 처리 \\
설화 가공 & Lv1 후크 / Lv2 장소-설화 연결 / Lv3 풀스크린 스토리(영구 캐시) \\
탐험 모드 & CoreLocation 지오펜스(도보 100m/차량 300m, 30초 체류) \\
지도 & Google Maps SDK + 웹 URL 폴백 \\
동행자 & 5종 페르소나(자동 배정+수동 변경), SSE 스트리밍 RAG \\
TTS & 페르소나별 음성 + 도착 인사·채팅·일지 낭독 3 포인트 \\
리뷰 & 5태그 + device\_id 익명, 자동 바텀시트, 일지 프롬프트 주입 \\
일지·이미지 & GPT-4o 단일 호출 JSON 출력 → 일지 + gpt-image-1 민화 \\
세션 복원 & SessionRestoreView 자동 감지·복원 \\
완료 동선 & ``내 코스'' 탭 자동 이동 + 1.5초 하이라이트 \\
공유 & iOS ShareLink + UIActivityViewController \\
\bottomrule
\end{tabularx}
\end{table}

\textbf{졸업 발표 시연 플로우 10단계.}\quad 본 시스템은 다음의 단일 여정을
끊김 없이 보여줄 수 있다:
(1) 취향 선택,
(2) AI 코스 추천 수신,
(3) 코스 미리보기 → 탐험 시작,
(4) 지도 화면(이동 경로 점선·다음 장소 카드),
(5) GPS 도착 시뮬레이션,
(6) AI 동행자 오버레이 등장 → 탭,
(7) 페르소나 채팅 2--3턴(텍스트+TTS),
(8) 장소 이동 반복(최소 2곳),
(9) 탐험 마치기 → Day 요약,
(10) 일지·민화 이미지 자동 생성.

\textbf{정량 지표.}\quad 표 6에 주요 정량 지표를 요약하였다.

\begin{table}[H]
\centering
{\small 표 6~~데이터·매핑 정량 지표}
\label{tab:metric}
\vspace{2pt}

\footnotesize
\begin{tabularx}{\columnwidth}{Xr}
\toprule
지표 & 값 \\
\midrule
큐레이션 코스 & 9{,}134개 \\
코스-장소 레코드 & 146{,}508개 \\
비짓제주 장소 (분류 전) & 3{,}284개 \\
설화 매핑 대상 관광지 (분류 후) & \textbf{820개} \\
매핑 보유 관광지 (524 / 820) & \textbf{64.0\%} \\
설화·민담 원본 & 505건 \\
ChromaDB 벡터 청크 (전체 적재) & 1{,}749개 \\
비짓제주 관광지 매핑 설화·민담 & \textbf{397 / 505 ($\sim$78\%)} \\
장소-설화 사전 매핑 (필터 전 → 후) & 134k → \textbf{30{,}947} ($-77\%$) \\
코스 추천에 사용된 매핑 & \textbf{6{,}726} (gps\_assist 제외) \\
일지+이미지 평균 응답 시간 & 25--30초 \\
\bottomrule
\end{tabularx}
\end{table}

% ── 6. 토의 ─────────────────────────────────────────────
\section{토~~~의}

\subsection{설계 의사결정의 Trade-off}

\textbf{왜 코스→설화 매핑인가.}\quad 설화→코스 역설계는 ``설화가 여행의
이유''라는 컨셉적 매력이 있으나, 설화 GPS 분포가 희박한 지역에서는 코스가
지리적으로 비현실적이 된다. 본 연구는 \emph{실제 관광객이 검증한 코스의
현실성}을 1순위로 두고, 그 위에 설화를 레이어로 얹는 방향을 선택하였다.
이는 ``AI 자유 생성의 자유도''와 ``실제 검증 데이터의 신뢰성'' 사이의 명시적
trade-off이며, 모바일 관광 앱에서는 후자가 더 결정적이라고 판단하였다.

\textbf{왜 페르소나 분리인가.}\quad 단일 챗봇으로 모든 설화 유형을 다루면
프롬프트가 짧고 유지보수가 단순하지만, 캐릭터 일관성이 ``유용한 비서''
어조로 수렴하여 무속신화·해녀·도깨비 등 서로 다른 세계관을 동일 어조로
전달하게 된다. 본 연구는 \emph{설화 유형별 5종 페르소나}로 분리하여 캐릭터
일관성과 적합성을 확보하였고, 자동 배정 규칙으로 사용자 선택 부담을 제거하면서
수동 변경 UI로 자유도도 함께 제공하였다.

\textbf{왜 영구 캐시인가.}\quad 설화 3단계 가공의 Lv2·Lv3는 LLM 호출이
필요하지만, 동일 (장소, 설화) 조합은 \emph{항상 동일 결과를 가져야 일관성
있는 여행 경험}이 된다. 매 요청마다 LLM을 호출하면 비용·지연뿐 아니라
동일 카드에서 매번 다른 문구가 나오는 일관성 문제가 발생한다. 영구 캐시
테이블 2종으로 ``한 번 생성, 영구 재사용'' 패턴을 적용하여 사용자 경험과
운영 비용을 동시에 최적화하였다.

\textbf{왜 핵심 기능에 ReAct 에이전트를 쓰지 않았는가.}\quad 코스 추천·동행자
채팅·일지 생성 등 핵심 경험은 응답 지연이 사용자 이탈에 직결되므로 도구
호출 라운드를 누적시키는 ReAct 패턴 대신 \emph{결정론적 함수와 단일 LLM
호출의 조합}을 채택하였다. ReAct는 매 추론 라운드마다 LLM을 재호출하여
지연이 누적되며, SSE 스트리밍과의 결합도 복잡해진다. 동행자 채팅은 사용자가
``지금 이 장소의 이야기''를 즉시 듣고 싶어 하는 맥락이므로, 도구 자율
호출의 자유도보다 한 번의 호출에 RAG 컨텍스트와 페르소나 프롬프트를 모두
주입하는 결정론적 구조가 우월하다. 본 시스템은 비용·지연·환각 통제를 위해
LLM 호출을 동행자 채팅, 일지·이미지 생성, 설화 3단계 가공이라는 세 지점에
한정하고, 각 호출 모두 ReAct 없이 단일 호출로 완결되도록 설계하였다.

\subsection{한계}

본 연구는 다음 한계를 가진다. 첫째, 공공데이터 설화·민담 505건 중 비짓제주
관광지와 매핑되지 않은 108건은 ChromaDB에 적재되어 있으나 동행자 채팅 검색
대상에서 제외된다. 본문이 매우 짧거나 매칭 가능한 비짓제주 관광지 지명이
본문에 등장하지 않는 경우로, 향후 자유 발화 챗봇이나 별도 콘텐츠 카드 등
다른 활용처를 검토할 수 있다. 둘째, 5종 페르소나에 대한 사용자 평가 데이터가
부재하여, 어느 페르소나가 어떤 사용자 군에 가장 효과적인지에 대한 정량적
근거가 아직 없다. 셋째, 실기기 현장 GPS 테스트가 시뮬레이션 위주로 진행되어,
제주 현지의 다중 경로·도심 협곡 환경에서의 지오펜스 신뢰성 검증이 충분하지
않다.

% ── 7. 결론 및 향후 과제 ─────────────────────────────
\section{결론 및 향후 과제}

본 연구는 ``실제 검증된 여행 코스 위에 제주 설화를 얹는다''는 컨셉으로,
계획·탐험·기록의 단일 여정을 책임지는 iOS 모바일 서비스를 설계·구현하였다.
LangGraph 인터페이스를 모방한 2단계 코스 추천 함수 파이프라인, 5종 페르소나 LLM
동행자와 페르소나 전용 TTS, 설화 3단계 가공과 영구 캐시, KTO+키워드 3단계
장소 분류 필터, GPT-4o 단일 호출의 일지·민화 동시 생성, 세션 복원,
``내 코스'' 자동 이동 동선 등 핵심 기능 파이프라인이 모두 완성되었다. 데이터 측면에서는 비짓제주
3{,}284개 장소를 820개 관광지로 한정하고 그중 524개에 설화·민담을 매핑하여 사전 매핑 테이블의
노이즈를 77\% 제거(134k → 30{,}947)하였으며, 코스 추천에는 GPS 보조 매핑을 제외한 6{,}726건만
사용한다. 설화·민담 505건은 ChromaDB에 1{,}749 청크로 적재하되, 비짓제주 관광지와 직접
매핑된 397 설화·민담의 청크만 동행자 채팅의 의미 검색 대상으로 한정한다.

\textbf{향후 과제.}\quad 본 서비스는 졸업 발표 이후 App Store 출시를 통해
실사용자 검증 단계로 진입한다. 시연 데모는 mock data 한계가 있으나, 실제
사용자 사용 데이터로 가설을 검증하는 단계가 학문적·실무적 신뢰도를 결정한다.
구체적으로 (1) TestFlight 베타와 매주 사용자 인터뷰로 페르소나 적합성·일지
만족도를 정량 평가하고, (2) 2026 관광데이터 활용 공모전(예비 합격, 10월
최종 심사)에 출품하며, (3) 외국인 관광객 대상 다국어 지원(KTO 영문
관광정보 API 연동), (4) 사용자 행동 로그 기반 페르소나·코스 개인화,
(5) 현장 GPS 정밀도 검증을 추진할 계획이다.

% ── 참고문헌 ────────────────────────────────────────────
\section*{참 고 문 헌}

\begin{small}
\begin{enumerate}
\setlength{\itemsep}{1pt}
\item P. Lewis et al., ``Retrieval-Augmented Generation for
  Knowledge-Intensive NLP Tasks,'' \textit{NeurIPS}, 2020.
\item 제주특별자치도, 제주 설화정보 공공데이터 API (E07), 2023.
\item 제주특별자치도, 제주 민담정보 공공데이터 API (E08), 2023.
\item 제주관광공사, 제주관광정보시스템(VISIT JEJU) 여행장소 및
  여행세부일정 공공데이터, 2026.
\item Chroma, ``Chroma: the AI-native open-source embedding database,''
  2023. [Online]. Available: \textit{https://www.trychroma.com}
\item OpenAI, ``Text Embedding Models,''
  \textit{platform.openai.com/docs/guides/embeddings}, 2024.
\item Google, ``Google Maps Platform -- Geocoding API \& Directions
  API, iOS SDK,'' \textit{developers.google.com/maps}, 2024.
\item S. Yao et al., ``ReAct: Synergizing Reasoning and Acting in
  Language Models,'' \textit{ICLR}, 2023.
\item LangChain, ``LangGraph: Build Stateful Multi-Actor
  Applications,'' \textit{langchain-ai.github.io/langgraph}, 2024.
\item 한국관광공사(KTO), TourAPI 공공데이터 서비스, 2024.
\item J. S. Park et al., ``Generative Agents: Interactive Simulacra of
  Human Behavior,'' \textit{Proc. UIST}, 2023.
\item OpenAI, ``GPT-4o System Card,''
  \textit{openai.com/index/gpt-4o-system-card}, 2024.
\item OpenAI, ``Image Generation (gpt-image-1 / DALL$\cdot$E),''
  \textit{platform.openai.com/docs/guides/images}, 2024.
\item W3C, ``Server-Sent Events,''
  \textit{html.spec.whatwg.org/multipage/server-sent-events.html}, 2024.
\end{enumerate}
\end{small}

\end{document}
